\section{Introduction}
How small a container holds $n$ copies of a given shape? For unit squares in a square the question goes back to \citet{erdosgraham}; catalogues now cover squares, triangles, hexagons and cubes in squares, triangles, circles and cubes \citep{friedmancenter,squaresproject}, and most entries are best known packings found by search rather than proved optimal. Most searches move rigid pieces: simulated annealing \citep{kirkpatrick}, perturbation and compression \citep{Gensane2004,gensane}, growing the pieces in size \citep{lubachevsky}, overlap minimisation in a shrinking container \citep{Umetani2009,Gardeyn2025}, or global solvers on exact non-overlap models \citep{Berthold2026}.

This paper changes the pieces instead. \emph{Hardening} starts every piece as the disk inscribed in it, compresses the disks under inward pressure on the container, and deforms them into the target shape while the pressure continues. It is graduated non-convexity \citep{blake} applied to geometry: a disk has no orientation, so the pieces are already densely packed when their corners appear and rotation becomes a choice. The shape family it moves along, a polygon dilated by a disk at fixed inradius, is not new: it is the spheropolygon of discrete-element modelling \citep{AlonsoMarroquin2008}, and \citet{Jin2015} moved dense periodic packings along its 3D analogue in both directions between tetrahedra and spheres. Our contribution is a controlled finite-container comparison of this optimisation path toward a prescribed polygon, with explicit tuning, budget and postprocessing limits (Section~\ref{sec:related}); it is not a claim to have invented shape evolution. On squares it recovers several classical packings from random starts, and the earlier companion project reported a twelve-cube feasible packing \citep{softtorigidpacking}; that prior discovery is not a new record of the controlled study here.

Here we ask whether hardening is a search worth running alongside existing ones. We compare three paths through shape space (hardening, growing, and rigid pieces on the same schedule) and two classical baselines on seven families of shapes and containers (\totalruns{} stored run records in \ncampaigns{} search campaigns, including deterministic C16 reuse rather than that many independent trials, plus one tightening-only campaign), with the explicit step and pair-evaluation budget design of Section~\ref{sec:method} and comparable grid-search effort. Since a packing search is judged by the best packing it finds, the comparison is by the \emph{lowest} container size each method reaches in a fixed number of runs, after constrained numerical tightening; the procedure being compared is the batch of runs, not a single run \citep{Birattari2005}. Our findings:
\begin{itemize}\itemsep1pt
\item Hardening is not a better default. Untuned, it looked better than rigid starts; tuned with comparable effort, rigid starts reach the lowest size most often over the seven families (Table~\ref{tab:main}).
\item It is a different search. On squares in a square it reaches the lowest size on \CsevenSquInSquHardenLowest{} of \CsevenSquInSquN{} instances and is the only method that low on \CsevenSquInSquHardenSole, including Trump's $n=11$ and Wainwright's $n=19$ packings, which rigid starts reach in none of our campaigns. In a targeted follow-up campaign at ten times the budget it is the only one of five methods to reach the best known packings for \LongHardenOnlyReachedNewNs; it also does so at $n=19$ and 29, but those hardening and rigid runs repeat an earlier campaign (Section~\ref{sec:learned}).
\item Exploratory ablations show that neither the disk-compressed start nor matched area growth reproduces hardening's pattern (Table~\ref{tab:ablate}), consistent with an effect of the gradual path; the pre-registered tests on squares in a square are inconclusive, and hardening's unique packings end with more tilted pieces (Figure~\ref{fig:mechanism}) without that isolating orientation as the cause.
\item The bias has a cost. Hardening is worse on triangles and on four families withheld from tuning and method analysis, apart from the disclosed preregistration geometry smoke tests, and the pilot policy we tested, which chooses the path per instance from 8 runs of each, showed no detected advantage over always starting rigid (Table~\ref{tab:decision}). An earlier indication that it gains more from longer runs did not replicate.
\end{itemize}
Search campaigns were planned in writing before they ran; amendments, exploratory analyses and post-hoc checks are identified separately. Raw outputs and saved analyses are versioned; selected best/hit trajectories have decimated replays. The original event export covers C00--C10, with later analyses also consuming campaign files directly. Online Resource 1 maps the evidence and distinguishes saved-data inspection from rerunning optimization. Section~\ref{sec:matters} gives the practical rule that follows and describes downstream reuse of the earlier release.
